% RSS 2027 Stage 1 extended abstract (due 2026-12-04 AoE). Limit: 5 pages of content + 1 page
% of references. Anonymized: no author names, affiliations, acknowledgements or links.
% Preamble follows the RSS LaTeX template (IEEEtran conference, Times, natbib numbers,
% plainnat). The official template (roboticsconference.org/docs/paper-template-latex.tar.gz)
% could not be downloaded from the drafting environment: before submission, diff this
% preamble against the official root.tex and use its .cls/.bst if they differ.
% Compile with pdflatex (Type-1 fonts): pdflatex main; bibtex main; pdflatex main; pdflatex main
\documentclass[conference]{IEEEtran}
\usepackage{times}
\usepackage[numbers]{natbib}
\usepackage{multicol}
\usepackage[bookmarks=true]{hyperref}

\usepackage[T1]{fontenc}
\usepackage{amsmath,amssymb}
\usepackage{booktabs}
\usepackage{array}
\usepackage{graphicx}
\usepackage{xcolor}
\usepackage{tikz}
\usetikzlibrary{arrows.meta,positioning}
\usepackage{pgfplots}
\pgfplotsset{compat=1.17}

\newtheorem{proposition}{Proposition}

% Draft notes. Set \notesfalse for the submission build: every note disappears.
\newif\ifnotes \notestrue
\newcommand{\pending}[1]{\ifnotes\textcolor{red}{[PENDING: #1]}\fi}
\newcommand{\verifynote}[1]{\ifnotes\textcolor{orange}{[CHECK: #1]}\fi}
\newcommand{\todo}[1]{\ifnotes\textcolor{orange}{[TODO: #1]}\fi}

\pdfinfo{
   /Author ()
   /Title  (Time Is the Baseline to Beat: Lead Time over a Timer as the Test for VLA Failure Monitors)
   /Subject (Robot learning; failure detection)
   /Keywords (vision-language-action;failure detection;conformal calibration;runtime monitoring)
}

\begin{document}

\title{Time Is the Baseline to Beat:\\Lead Time over a Timer as the Test for VLA Failure Monitors}

\author{Author Names Omitted for Anonymous Review. Paper-ID [add your ID here]}

\maketitle

\begin{abstract}
On RoboCasa365, a failed vision-language-action (VLA) rollout is one that reaches its time
limit. Calibrated at a 5\% false-positive rate (FPR), a timer flags every one of 380
prospective failures, so recall cannot separate a learned monitor from elapsed time; only
warning lead can. We alarm when \emph{either} a policy-internal score or a timer crosses its
threshold, calibrate the pair with a finite-sample FPR guarantee, and compare with a per-task
timer P built from each task's success durations. On one VLA and 38 tasks, P nearly matches
the SAFE failure score fused with a single timer, and on a fresh preregistered cohort of 760
rollouts SAFE adds 0.030 of the time limit to P (95\% CI 0.007--0.093), below the margin of
0.05 we declared in advance. A sequence detector trained for early separation on the same
features, Hide-and-Seek, passes the same registered test: it adds 0.071 (CI 0.022--0.157) at
0.061 FPR. Retrained on zero inputs, the same model matches its ROC-AUC but adds no lead; given
only each rollout's first inference, it adds none either. Elapsed time explains how well
monitors separate failures; the extra warning develops during the rollout. Stage~2 tests a
second VLA and whether earlier alarms raise success when they trigger a retry.
\end{abstract}

\section{Introduction}

Runtime monitors for VLA policies read the policy's hidden state and predict eventual
failure~\cite{gu2025safe}, or score how unusual that state is relative to successful
rollouts~\cite{xu2025faildetect,romer2025fiper}. They are compared with each other, but not
with what the robot knows for free: how long the attempt has been running.

In RoboCasa365, as in most simulated manipulation benchmarks, an episode ends at first
success or at a time limit, and in our test cohort every one of 380 failed rollouts ran to
its limit. A timer at 90\% of the limit therefore flags all of them at 4.7\% FPR. Against it,
a monitor can add only one thing: an earlier alarm at the same false-alarm rate, which leaves
time to stop, reset or ask for help. We make this lead time the test, with three
contributions.
\begin{enumerate}
  \item \textbf{A testable hypothesis and endpoint.} We define the warning lead of an alarm
  as the fraction of the time limit left when it fires (zero for a missed failure) and state
  four hypotheses with decision rules fixed before the confirmatory data were analysed
  (Sec.~\ref{sec:hypotheses}).
  \item \textbf{Calibrated OR-fusion with a guarantee.} We alarm when either a
  policy-internal score or the timer crosses its threshold, and calibrate both thresholds
  along a nested path with a split-conformal rule that bounds the FPR on a new successful
  rollout by $\alpha$ in finite samples (Proposition~\ref{prop:fpr}).
  \item \textbf{Exploratory and confirmatory evidence.} A per-task timer nearly matches SAFE
  fused with a single timer, and most failures are still recoverable at the limit
  (Sec.~\ref{sec:evidence}). On fresh preregistered cohorts SAFE adds 0.030 of the limit over
  the per-task timer, below our 0.05 margin, while Hide-and-Seek~\cite{park2026hideseek} adds
  0.071, and controls show that its extra warning is not learned elapsed time
  (Sec.~\ref{sec:confirm}).
\end{enumerate}
Elapsed time explains most of how well monitors \emph{separate} failures; lead over a
per-task timer shows which ones also warn \emph{earlier}. On timeout-defined benchmarks,
monitors should be compared with a per-task timer and evaluated on failures that end early.

% Related Work for the RSS 2027 Stage 1 draft.
% Merged from /mnt/project-files/paper/related_work.tex (related-work thread, 2026-10-17),
% shortened for the 5-page limit and reframed around the timer and lead time.
% Claims about other papers were checked by that thread against abstracts or proceedings
% pages; RESCUE was read from its arXiv abstract only.

\section{Related Work}
\label{sec:related}

\textbf{Failure detection from policy internals.}
SAFE learns a failure score from a VLA's last-layer features and thresholds it with a
functional conformal band~\cite{gu2025safe}; we reuse its implementation. Follow-ups change
the training data~\cite{rajaprakash2026safecast}, labels~\cite{park2026hideseek,ma2026failurespot},
features~\cite{yang2026triinfo,zhu2026vlascope} or policy~\cite{soddy2026memoryvla}. None
reports an elapsed-time detector. SAFE removes the duration shortcut by truncating rollouts to
equal length~\cite{gu2025safe}, which a deployed monitor cannot do. Duration is a shortcut
feature~\cite{geirhos2020shortcut}, but its causal form is legitimate evidence online, so we
make it the baseline.

\textbf{Action-space and success-only detectors.}
STAC~\cite{agia2024sentinel} and VLA-FAIL~\cite{seligmann2026vlafail} compare overlapping
action chunks, and ActProbe~\cite{huang2026actprobe} feeds chunk consistency and magnitude to
a task-conditioned LSTM, beating internal-feature baselines on RoboCasa. Our policy executes
each chunk in full, so overlap signals are undefined; chunk magnitude and a second flow sample
are Stage~2 baselines. A task-conditioned recurrent detector can also learn how long each task
takes, which our timers make explicit. FAIL-Detect~\cite{xu2025faildetect} and
FIPER~\cite{romer2025fiper} score inputs, outputs or embeddings against successful rollouts
only. FIPER alarms when two conformal scores \emph{both} exceed their thresholds, which
suppresses false alarms; we alarm when \emph{either} crosses, because the policy score and
the timer catch failures at different times, and calibrate the pair jointly.

\textbf{World-model verifiers and VLM monitors.}
CheckVLA~\cite{liu2026checkvla} is the closest work on our benchmark. It checks each action
chunk against an action-conditioned world model on RoboCasa365, calibrates a conformal
threshold to a 5\% episode-level rate of unnecessary interventions, and repairs the chunk
with the same VLA. Foresight~\cite{wang2026foresight} uses world-model latents. Neither compares against
elapsed time, and either risk could be fused with a timer as we fuse SAFE. LLM and VLM
monitors~\cite{agia2024sentinel,liu2023reflect,sinha2024realtime,zhou2025codeasmonitor,duan2025aha}
are slower and complementary.

\textbf{Time, calibration and intervention.}
VLA confidence drifts over a task~\cite{zollo2026calibration}. Our time-risk baseline is a per-task
survival-style estimate~\cite{kaplan1958,cox1972}, and lead time at a fixed false-alarm
budget is the accuracy--earliness trade-off of early time-series
classification~\cite{xing2012early}. Split-conformal thresholds give finite-sample
false-positive control under exchangeability~\cite{angelopoulos2021gentle}; we extend the
construction to a two-signal OR rule (Sec.~\ref{sec:guarantee}). RESCUE~\cite{ahmedt2026rescue}
finds that good detection scores do not guarantee lower intervention cost, and CheckVLA and
FLARE~\cite{zhao2026flare} couple detection to repair. Hence our closed-loop Stage~2 test: an
earlier alarm matters only if acting on it raises success.


\section{Problem Setting}
\label{sec:setting}

\textbf{Time.} The simulator runs \emph{control steps} at RoboCasa's
default 20\,Hz. The policy executes each 16-step action chunk in full, so \emph{policy
inferences} $t = 1, 2, \dots$ (0.8\,s each) are our clock, and a monitor decides only at an
inference. Task $k$'s limit of $T_k$ control steps is $H_k = \lceil T_k/16 \rceil$ inferences
(29 to 272, or 22.5 to 217.5\,s). Every fraction of time ($t/H_k$, $c$, $\ell$) is a fraction
of $H_k$.


\textbf{Monitoring.} A monitor watches a rollout of a known task $k$ and decides at each
inference $t$ whether to alarm, using only the policy features so far, $t$, the task and
$H_k$. The first crossing is the alarm $T$ (in inferences). FPR is the fraction of successful
rollouts with any alarm. A rollout of length $L$ inferences ends at success or at $H_k$; for
a failure ($L = H_k$) the \emph{warning lead} is $\ell = (L - T)/H_k$, or 0 if missed, and an
alarm at the final inference counts with $\ell = 0$. The mean lead $\bar\ell$ combines recall
and earliness. Detectors are compared at operating points fixed on a calibration cohort under
FPR $\le \alpha = 0.05$. Confidence intervals come from a paired bootstrap that resamples
tasks, then rollouts within each task (2,000 replicates).

\textbf{Scope.} Whether elapsed time is a strong baseline depends on the termination rule.
Where successes stop at completion (RoboCasa, FIPER's tasks, deployment), it is. Where every
rollout runs to one fixed horizon, as in SAFE's real-robot set (in its release, 98\% of
failures and $\ge 30$ successes per task stop at one step) or SimplerEnv, it is uninformative,
and early-warning claims need failure-time annotations, which SAFE did not release.

\section{Hypotheses}
\label{sec:hypotheses}

We test four hypotheses. H1 and H2 were evaluated on a fresh confirmatory cohort with the
decision rules below, registered as an amendment before any timer or lead quantity was
computed on that cohort, and registered again for a second score before that score was
computed on it (Sec.~\ref{sec:confirm}); H3 and H4 are the Stage~2 extensions
(Sec.~\ref{sec:plan}).

\textbf{H1 (lead time, primary).} At thresholds calibrated for FPR $\le 0.05$, OR-fusing a
policy-internal score with the per-task timer (Sec.~\ref{sec:method}), the strongest time-only
monitor we know, increases the mean warning lead over that timer alone.
\emph{Supported} if the paired bootstrap 95\% CI of $\Delta\bar\ell$ lies above zero, the
point gain is at least $0.05$ of the time limit, and the fused test FPR is at most $0.075$
with a Wilson 95\% upper bound at most $0.10$. As a secondary check, the fused detector must
also beat a timer whose fraction is tuned \emph{on the test cohort} to the fused detector's
test FPR, which favours the timer.

\textbf{H2 (the gain needs policy evidence).} Fusions that contain no policy-internal
evidence (Gaussian noise, or SAFE scores resampled within task, each fused with the per-task
timer) gain less than $0.05$ of the time limit over the per-task timer.

\textbf{H3 (generality).} H1 holds for a second VLA with a different architecture.%
\todo{Must change to match E2, which drops the 0.05 margin; author's call. Proposal: on
RLDX-1's 20 tasks, (a) every failure runs to the time limit, (b) P leads the global timer, and
(c) SAFE + P gains lead over P with a within-task 95\% CI above zero, under the FPR rule of
Sec.~\ref{sec:plan}.}

\textbf{H4 (utility).} When an alarm triggers a reset and retry, the OR-fused monitor yields
a higher task success rate than the timer within the same per-task time budget.%
\todo{Must change to match E3; author's call. Proposal: within a budget of $1.25\,H_k$,
reset and retry triggered by P yields a higher success rate than no monitor; SAFE + P against
P is descriptive.}

Only H1 and H2 are confirmatory; every other interval in this paper is descriptive. We
expected H1 to fail for SAFE (opened-data gain 0.007); it did, with a gain of 0.030 whose CI
excludes zero. A second amendment then applied the same rules to Hide-and-Seek, which gained
0.083 at matched FPR on the opened cohort; it passed (Sec.~\ref{sec:confirm}).

\section{Method}
\label{sec:method}


\subsection{Signals}

\textbf{Policy score.} At each inference we take the final-layer action-token hidden state
of the policy's action head (1,024 dimensions). The SAFE MLP~\cite{gu2025safe} maps it to a
failure score $u^{(j)}_t$ for each of three training seeds $j$. With per-task location
$m^{(j)}_k$ and scale $\sigma^{(j)}_k$ estimated on development rollouts, the policy score is
\begin{equation}
  s_t = \frac{1}{3}\sum_{j=1}^{3} \frac{\max_{t' \le t} u^{(j)}_{t'} - m^{(j)}_k}{\sigma^{(j)}_k}.
\end{equation}
The running maximum makes $s_t$ monotone and does not change when a threshold is first
crossed. Success-only alternatives (Mahalanobis~\cite{lee2018mahalanobis},
kNN~\cite{sun2022knn}, PCA-KMeans and RND~\cite{burda2019rnd}) are fitted on successful
development rollouts and aggregated the same way.

\textbf{Sequence score and controls.} Hide-and-Seek~\cite{park2026hideseek} trains an LSTM on
a 128-dimensional PCA of the same hidden state with rollout labels only: a margin loss pushes
each failure's peak score above every success's, and an onset loss acts within failures. Its
scores are normalized and aggregated like SAFE's. A recurrent model can learn how long each
task takes, so we retrain it unchanged on (i) zero inputs, (ii) the features of a random
same-task training rollout at the same inference, and (iii) the rollout's own first-inference
features, repeated. They see (i) elapsed time, (ii) also the task's typical trajectory, and
(iii) also the episode's initial appearance. ActProbe~\cite{huang2026actprobe} reduces to chunk
magnitude here, as our chunks never overlap.

\textbf{Timers.} The global timer alarms once $t \ge c\,H_k$ for a fraction $c \in (0, 1]$
shared by all tasks; it uses only the task's time limit. The per-task timer uses
$c_k = \min(1, \rho_k + \delta)$, where $\rho_k$ is the 95th percentile of task $k$'s
development success durations as fractions of $H_k$ and the shared margin $\delta$ is set by
Proposition~\ref{prop:fpr} on calibration successes. It needs successful rollouts of each
task, but no failures and no policy features. We fixed the 95th percentile after auditing the
opened cohort, so we report 0.75, 0.90 and 1.0 (the longest success) as sensitivity.

\textbf{Time risk (baseline).} The survival-style alternative estimates
$r_t = \Pr(\text{failure} \mid k, \text{still running at } t)$ from development rollouts: with
$n_t$ of them still running at $t$ and $f_t$ of those failing, it fits an increasing
isotonic regression in $t/H_k$ to $(f_t + 0.5)/(n_t + 1)$, weighted by $n_t$. Because every
failure runs to $H_k$, $r_t$ is a task-specific timer; its smoothing keeps it below a high
threshold in tasks with few development failures, which is why it misses failures the
plain timer catches.

\subsection{OR-fusion}

The fused monitor alarms at
\begin{equation}
  T = \min\{\, t : s_t > \tau_s \ \text{or}\ t \ge c\,H_k \,\}.
\end{equation}
Neither signal is restricted in time: the policy score may alarm early, and the timer
guarantees an alarm on every failure, at $c\,H_k$ at the latest. The lead gain over the
timer comes only from failures on which the policy score crosses first.

\subsection{Calibration with a finite-sample FPR guarantee}
\label{sec:guarantee}

A successful rollout $i$ of length $L_i$ (in inferences) raises a false alarm under
$(\tau_s, c)$ if and only if $M_i = \max_{t \le L_i} s_t > \tau_s$ or $\rho_i = L_i/H_{k_i}
\ge c$, so its false alarm depends only on the pair $(M_i, \rho_i)$. Fix in advance a path
$\gamma(\lambda) = (\tau_s(\lambda), c(\lambda))$, $\lambda \in [0, 1]$, with both
coordinates continuous and nondecreasing, so that a larger $\lambda$ is more conservative and
$\gamma(1)$ never alarms on a success. For each calibration success define
\begin{equation}
  V_i = \min\{\lambda : M_i \le \tau_s(\lambda) \ \text{and}\ \rho_i < c(\lambda)\},
\end{equation}
the least conservative point on the path at which rollout $i$ stays silent. Because the path
is monotone, rollout $i$ is silent at $\lambda$ exactly when $\lambda \ge V_i$.\footnote{The
strict $\rho_i < c(\lambda)$ makes the set right-open; taking $c$ on a finite grid, as the
timer's resolution of one inference does anyway, makes the minimum attained.}

\begin{proposition}
\label{prop:fpr}
Let $n$ calibration successes and a new successful rollout be exchangeable, and let the path
$\gamma$ not depend on the calibration successes. Set $\hat\lambda = V_{(\lceil (n+1)(1-\alpha)
\rceil)}$, the $\lceil (n+1)(1-\alpha) \rceil$-th smallest $V_i$ (or $\hat\lambda = 1$ if
that rank exceeds $n$). Then the OR-fused monitor at $\gamma(\hat\lambda)$ raises a false
alarm on the new success with probability at most $\alpha$.
\end{proposition}
\begin{IEEEproof}
The new rollout alarms at $\hat\lambda$ if and only if $V_{n+1} > \hat\lambda = V_{(r)}$ with
$r = \lceil (n+1)(1-\alpha) \rceil$. By exchangeability, the rank of $V_{n+1}$ among
$V_1, \dots, V_{n+1}$ is uniform on $\{1, \dots, n+1\}$ (ties broken at random), so
$\Pr(V_{n+1} > V_{(r)}) \le 1 - r/(n+1) \le \alpha$; this is the split-conformal quantile
lemma~\cite{angelopoulos2021gentle} applied to the scalar $V$.
\end{IEEEproof}

The path may use any data except calibration successes, such as development rollouts or
calibration \emph{failures}. The guarantee is marginal over tasks; a per-task version needs
19 calibration successes per task.

The preliminary results use the original protocol's joint search: among all pairs of
calibration maxima with calibration FPR $\le \alpha$, keep the one with the largest mean lead
(highest recall for time-risk fusions). This best-of-many selection on 114 successes has no
guarantee and is optimistic (test FPR 0.058--0.068 against 0.044 on calibration). The
per-task timer rows, and the confirmatory test, use Proposition~\ref{prop:fpr} with the path
$\delta = \lambda$, $\tau_s = u + s\lambda$ and $(u, s)$ chosen on calibration failures.

\section{Preliminary Evidence}
\label{sec:evidence}

% Qualitative figure: one failed OpenCabinet rollout (prospective test), frames + SAFE score.
% Data: rollout_curves.csv and frames.json from make_qualitative_data.py.
% Frames: put f1.png..f4.png in \qualdir/frames/ (see README.md); grey placeholders until then.
% Needs: tikz, pgfplots, graphicx. Set \qualdir before \input if the folder is elsewhere.
\providecommand{\qualdir}{figures/qualitative}
\newcommand{\qualframe}[2]{%
  \begin{minipage}[t]{0.235\columnwidth}\centering
  \IfFileExists{\qualdir/frames/#1.png}%
    {\includegraphics[width=\linewidth]{\qualdir/frames/#1.png}}%
    {\tikz\node[draw=gray!60, fill=gray!12, minimum width=\linewidth, minimum height=0.75\linewidth,
       font=\tiny, text=gray!70!black, inner sep=0]{frame #1};}\\[-1pt]
  {\scriptsize #2}\end{minipage}}
\begin{figure}[t]
\centering
\qualframe{f1}{$t=10$}\hfill
\qualframe{f2}{$t=30$: SAFE}\hfill
\qualframe{f3}{$t=60$: timer}\hfill
\qualframe{f4}{$t=66$: limit}\\[3pt]
\begin{tikzpicture}
\begin{axis}[width=\columnwidth, height=3.6cm, xmin=0, xmax=1, ymin=-1.9, ymax=3.1,
  xlabel={elapsed fraction of the time limit $t/H_k$}, ylabel={SAFE score $s_t$},
  label style={font=\scriptsize}, tick label style={font=\scriptsize},
  xlabel shift=-3pt, ylabel shift=-4pt,
  legend style={font=\tiny, at={(0.01,0.99)}, anchor=north west, draw=none, fill=white, fill opacity=0.8, text opacity=1},
  legend cell align=left, grid=major, grid style={gray!15}, unbounded coords=jump]
\addplot[blue!75!black, thick] table[col sep=comma, x=frac, y=fail_score] {\qualdir/rollout_curves.csv};
\addlegendentry{failed rollout}
\addplot[green!50!black, thick, densely dashed] table[col sep=comma, x=frac, y=succ_score] {\qualdir/rollout_curves.csv};
\addlegendentry{successful rollout, same task}
\addplot[red, dashed, domain=0:1, samples=2] {0.648};
\addlegendentry{SAFE threshold $\tau_s$}
\draw[gray!70!black, thick, dotted] (axis cs:0.9027,-1.9) -- (axis cs:0.9027,3.1)
  node[pos=0.93, left, font=\tiny, text=gray!40!black] {timer, $c=0.90$};
\draw[orange!80!black, thick, dashdotted] (axis cs:0.5758,-1.9) -- (axis cs:0.5758,3.1)
  node[pos=0.80, right, font=\tiny, text=orange!50!black] {per-task timer};
\fill[red!80!black] (axis cs:0.4545,0.648) circle (1.6pt);
\draw[<->, >=stealth, red!70!black] (axis cs:0.4545,-1.65) -- (axis cs:1,-1.65)
  node[midway, above, font=\tiny] {lead 36 inferences ($\approx$29\,s)};
\draw[<->, >=stealth, orange!70!black] (axis cs:0.5758,-0.85) -- (axis cs:1,-0.85)
  node[midway, above, font=\tiny] {28};
\draw[<->, >=stealth, gray!50!black] (axis cs:0.9091,0.0) -- (axis cs:1,0.0)
  node[midway, above, font=\tiny] {6};
\end{axis}
\end{tikzpicture}
\caption{One failed \emph{OpenCabinet} rollout from the prospective test (chosen to
illustrate; the median SAFE-first lead is smaller). The SAFE score crosses its calibrated
threshold at $t = 30$ of $H_k = 66$ inferences, 36 inferences before the time limit; the timer
alarms at $t = 60$, and the per-task timer, which never reads the policy, at $t = 38$. Dashed: the longest of the 9 (out of 10) successful test rollouts of this task that stay
below the threshold; it ends at $t = 33$.}
\label{fig:qualitative}
\end{figure}


\textbf{Setup.} We use the released Xiaomi-Robotics-1 RoboCasa365
checkpoint~\cite{xiaomi2026robotics}, which executes each 16-step action chunk in full, on the
38 official RoboCasa365 tasks~\cite{nasiriany2024robocasa,robocasa365} with published success
between 5 and 45 of 50 and enough failures (10 atomic, 28 composite; failures last 29--272
inferences). Disjoint cohorts serve one purpose each: development (1,824 rollouts) fits SAFE
and the time risk; calibration (3 successes and 3 failures per task; 454 rollouts including
those past the quota) sets thresholds; the prospective test (10 per class per task) is
evaluated once; a confirmatory cohort of the same size is used only in Sec.~\ref{sec:confirm}. SAFE, its
normalization and the time risk were frozen before the test was collected; the timers and
fused detectors were added later, with thresholds from calibration only, so the results below
are exploratory.

\begin{table}[t]
\centering
\caption{Prospective test (380 successes, 380 failures), thresholds from calibration at FPR
$\le 0.05$. $\bar\ell$: mean lead (miss = 0); Med: median lead (inferences) of detected
failures; @.75: failures alarmed by $0.75\,H_k$. $^\ast$Mean of 20 draws ($\bar\ell$
0.088--0.134). $^\dagger$Margin by Proposition~\ref{prop:fpr}.}
\label{tab:main}
\footnotesize
\setlength{\tabcolsep}{3pt}
\begin{tabular}{@{}lccccc@{}}
\toprule
Detector & Recall & FPR & $\bar\ell$ & Med & @.75 \\
\midrule
Timer ($c = 0.903$) & 1.000 & 0.047 & 0.090 & 11.5 & 0.000 \\
Time risk & 0.632 & 0.047 & 0.120 & 16.5 & 0.211 \\
Time risk + timer & 1.000 & 0.050 & 0.125 & 11.0 & 0.105 \\
Noise + timer$^\ast$ & 1.000 & 0.066 & 0.101 & -- & -- \\
SAFE only & 0.566 & 0.039 & 0.190 & 23.0 & 0.421 \\
SAFE + time risk & 0.955 & 0.058 & 0.184 & 12.0 & 0.311 \\
Per-task timer$^\dagger$ & 1.000 & 0.045 & 0.195 & 13.5 & 0.342 \\
SAFE + per-task timer$^\dagger$ & 1.000 & 0.042 & 0.202 & 14.0 & 0.389 \\
\textbf{SAFE + timer} & 1.000 & 0.058 & 0.207 & 10.0 & 0.421 \\
\bottomrule
\end{tabular}
\end{table}


\textbf{The timer is the baseline to beat.} The timer ($c = 0.903$) alarms on all 380 test
failures at 0.047 FPR with a fixed lead of 0.090 (Table~\ref{tab:main}). The time-risk curve
misses 140 failures, all in 14 tasks whose smoothed curve never reaches its threshold, so recall
gains measured against it (to 0.90--0.96 for SAFE and four success-only scores) vanish here.

\textbf{SAFE adds lead over a single timer.} SAFE + timer raises the mean lead to 0.207 (paired
gain 0.117, CI 0.060--0.178) and flags 42\% of failures by $0.75\,H_k$
(Fig.~\ref{fig:qualitative}), at 0.058 test FPR; noise fused with the timer reaches 0.101.

\textbf{A per-task timer removes the gain (H1, exploratory).} The per-task timer reaches
$\bar\ell = 0.195$ at 0.045 FPR without reading the policy. SAFE + timer leads it by an
insignificant 0.012 (CI $-0.044$ to $0.063$). SAFE + per-task timer reaches 0.202, a gain of 0.007 (CI $-0.003$ to
$0.019$), far below the 0.05 margin of H1.

\textbf{A sequence detector does warn earlier (exploratory).} Fused with P under
Proposition~\ref{prop:fpr}, Hide-and-Seek gains 0.098 over P (CI 0.037--0.158), and 0.083
(0.022--0.143) against P loosened to its own 21 test false alarms. In that matched comparison
SAFE gains 0.008, ActProbe 0.009, and the best of 100 noise draws 0.016. The zero-input
control matches Hide-and-Seek's ROC-AUC (0.888 against 0.878) and its recall when fused with
the time risk (0.974 against 0.976), yet adds 0.001 over P at matched FPR, and the same-task
shuffle $-0.004$: elapsed time explains separation, not earliness.

\textbf{No early warning, and mostly no early failure (exploratory).} No detector flags over
8\% of failures by half the limit. Is timeout-defined failure then trivial by construction? One
labeler, blind to alarms, marked the point of no return (after which the policy could not
finish in time without help) on 48 confirmatory-cohort failures (2 per atomic, 1 per composite
task): 38 (79\%) were still recoverable at the limit (atomic 17/20, composite 21/28), most
often slow or nearly successful. A second labeler agreed on 11 of 12 (Cohen's $\kappa = 0.80$).
Failure here mostly means not finishing in time, which is why elapsed time is so informative;
an early ``failure'' alarm mostly predicts slowness. On the 10 decided failures, P alarmed
after the point of no return on all, SAFE + P at or before it on 2 of 3 atomic and 0 of 7
composite ones (an illustration).

\textbf{A second policy, offline (RLDX-1, exploratory).} On existing RLDX-1 data (20 tasks,
100/100 test rollouts), all failures ran to the time limit, P reaches $\bar\ell = 0.271$
against 0.032 for the global timer, and SAFE (LSTM head) + P gains 0.034 over P (CI
0.008--0.067) at the same FPR, all on atomic tasks; the fused path was chosen in sample and
noise + P matched the gain in 2 of 20 draws at a higher FPR.

\textbf{The confound in a published benchmark (exploratory).} In FIPER's released
rollouts~\cite{romer2025fiper}, length alone separates failures with ROC-AUC 0.92--1.00 in
four of five tasks. A timer calibrated like FIPER's thresholds exceeds its published mean
accuracy (0.83 against 0.78) but alarms at 0.63 of the horizon against 0.30: FIPER's advantage
is earliness, the comparison we argue for (published against computed numbers, not a paired
re-run).

\section{Confirmatory Test of H1 and H2}
\label{sec:confirm}

\textbf{Protocol.} A fresh test cohort (new seeds, 10 per class per task) was collected with
SAFE frozen and preregistered for recall against the time risk. Before any timer or lead was
computed on it, but with its recall and FPR known, an amendment made the per-task timer P
($q = 0.95$) the comparator and $\Delta\bar\ell$ the endpoint, with the rules of
Sec.~\ref{sec:hypotheses}. A new calibration cohort of the same design (380 successes) sets
the thresholds of P and SAFE + P by Proposition~\ref{prop:fpr}. The decision interval
resamples both cohorts and refits thresholds in each of 2,000 replicates. Every test failure
again ran to $H_k$.

\begin{table}[t]
\centering
\caption{Confirmatory test (380 successes, 380 failures). P = per-task timer ($q = 0.95$).
Intervals from P down: registered decision bootstrap; above P: test only. $^\ast$Mean of 20
draws [maximum]; resampled: SAFE, H\&S. $^\dagger$Frozen original thresholds. H\&S:
Hide-and-Seek.}
\label{tab:confirm}
\footnotesize
\setlength{\tabcolsep}{3pt}
\begin{tabular}{@{}lcccl@{}}
\toprule
Detector & Recall & FPR & $\bar\ell$ & Gain vs P \\
\midrule
Global timer & 1.000 & 0.024 & 0.123 & $-0.102$ ($-0.147$, $-0.059$) \\
SAFE + time risk$^\dagger$ & 0.955 & 0.039 & 0.190 & $-0.035$ ($-0.066$, $-0.007$) \\
SAFE + global timer & 1.000 & 0.047 & 0.215 & $-0.011$ ($-0.068$, $0.045$) \\
P & 1.000 & 0.055 & 0.225 & -- \\
Noise + P$^\ast$ & -- & -- & -- & $0.002$ [$0.009$] \\
Resampled + P$^\ast$ & -- & -- & -- & $0.002$, $0.000$ [$0.007$] \\
\textbf{SAFE + P (H1)} & 1.000 & 0.061 & \textbf{0.255} & $\mathbf{0.030}$ ($0.007$, $0.093$) \\
First-inference + P & 1.000 & 0.055 & 0.226 & $0.001$ ($-0.015$, $0.036$) \\
Zero-input + P & 1.000 & 0.061 & 0.237 & $0.012$ ($0.001$, $0.044$) \\
Shuffle + P & 1.000 & 0.063 & 0.236 & $0.011$ ($0.000$, $0.064$) \\
\textbf{H\&S + P (Am.~2)} & 1.000 & 0.061 & \textbf{0.296} & $\mathbf{0.071}$ ($0.022$, $0.157$) \\
\bottomrule
\end{tabular}
\end{table}

\textbf{H1 fails for SAFE, but the gain is not zero.} SAFE + P leads P by 0.030 of the time
limit (Table~\ref{tab:confirm}). Its decision CI lies above zero and its test FPR (0.061,
Wilson 95\% upper bound 0.089) meets the FPR rules, but the point gain is below the 0.05
margin, so H1 fails.

\textbf{H2 is supported.} Noise and within-task resampled SAFE scores, fused with P and
calibrated the same way, gain at most 0.009 over P in 20 draws each. What SAFE adds is
policy evidence, not an artifact of fusing any score with a timer.

\textbf{Task duration carries most of SAFE's lead.} P alone leads the global timer by 0.102,
and the original SAFE + time-risk design, which beat the time risk on recall here, gives 0.035
\emph{less} lead than P. A per-task timer tuned on the test cohort to SAFE + P's FPR leaves
SAFE a gain of 0.015 (CI $-0.007$ to $0.041$). Descriptively, the four success-only scores
fused with P gain 0.020--0.034, and all five scores gain more on the 10 atomic tasks
(0.045--0.081) than on the 28 composite ones (0.008--0.027).

\textbf{Amendment 2: Hide-and-Seek passes H1 and H2.} It was registered before any
Hide-and-Seek score was computed on these cohorts. The retrained detector reproduced its
exploratory scores exactly, and the analysis reproduced SAFE + P before reporting.
Hide-and-Seek + P leads P by 0.071 (decision CI 0.022--0.157) at 0.061 FPR (Wilson bound
0.089), meeting every rule; noise and resampled scores gain at most 0.012. The three controls
gain 0.001--0.012. Hide-and-Seek beats the first-inference control by 0.070 (CI
0.014--0.153), so its warning develops during the rollout rather than reflecting a
hard-looking scene, and beats the zero-input and shuffle controls by 0.059 and 0.060 (CIs
above zero). Against SAFE + P its advantage, 0.041 (CI $-0.011$ to 0.103), is not significant.
The decision bootstrap skews upward (SAFE: 0.007--0.093 around 0.030), so here the point gain,
not the interval, carries the decision.

\section{Stage 2 Plan}
\label{sec:plan}

% Stage 2 sizing: docs/stage2_sizing.md, Results (branch claude/stage2-sizing, commit 173a0274).
% Offline: resample the opened cohorts, refit, count. Design decisions by the author, 2026-10-21.
\textbf{E2: Second policy (H3).} RLDX-1 on its 20 tasks (Sec.~\ref{sec:evidence}), fresh
calibration and test cohorts of 20 successes and 20 failures per task (about 2,500 rollouts),
with held-out calibration failures. Primary: SAFE + P against P, superiority (decision 95\% CI
above zero), conditional on these 20 tasks (rollouts resampled within task), under the FPR rule
below.\todo{Author's call: make Hide-and-Seek + P co-primary (an offline RLDX-1 check is
running), with the first-inference control as secondary.} Offline sizing gives power 0.92 at the offline gain of 0.034 (the CI excludes zero in every
simulated study, so the FPR rule decides); a 0.05 margin would have power $\le 0.09$, so we
drop it. Secondary, no power claim: the CI with tasks resampled (power
about 0.54). Descriptive stratum: the 14 atomic tasks. It replaces an atomic-only XR-1 test,
with at most 0.55 power at the 0.05 margin up to 60 per class per task and an optimistic
0.089 estimate (about 0.06 with all retained successes). With 5 held-out failures
per task, all sizing figures are optimistic.

\textbf{FPR rule.}\todo{Proposed rule; author to confirm.} Calibration keeps the conformal
target $\alpha = 0.05$. H1's fixed limit (FPR $\le 0.075$, Wilson bound $\le 0.10$) rejects
good detectors in small cohorts; instead, false alarms must not reject ``true FPR $\le
0.075$'' in a one-sided exact binomial test at level 0.05: at most the 95th percentile of
Binomial($n_s$, 0.075) for $n_s$ test successes (39 of 400 in E2). This is looser, so E2's power stays $\ge 0.92$.

\textbf{E3: Reset and retry (H4).} On XR-1's 38 tasks\pending{E3 policy not chosen; XR-1
assumed (RLDX-1: about 480 per arm)}, an alarm ends the attempt and resets the scene while a
budget of $1.25\,H_k$ remains. Registered: P against no monitor, paired by seed (about 650
episodes per arm). SAFE + P is descriptive: separating it from P needs thousands to over a
million episodes per arm (about 1,450 on XR-1 atomic tasks at $1.0\,H_k$, on an optimistic
bound). As most failures are still recoverable at the limit, a ``more time, no reset'' arm
(run to $1.25\,H_k$) is the natural comparator for reset.\pending{sample size of the no-reset arm}

\textbf{E4: Checks.} Action-only monitors (chunk magnitude, second-sample disagreement);
richer time baselines (end-effector stall, discrete-time hazard); OR vs AND on full curves;
SAFE faithfulness; precision at deployment prevalence; unseen tasks; FIPER's code re-run
with the timer. \textbf{E5: Failures decided early.} Most failures have no early point of no
return, so E5 targets the minority that do, with several labelers, a written protocol and
adjudication (our two labelers' timings differed by 0.24--0.34 $H_k$), and reports lead from
that moment.

\section{Limitations}

All results are simulated shadow evaluations (monitors never intervened), from one policy
prospectively; a second (RLDX-1) was checked only offline, on existing data with a small test
set. Timers and fused detectors were defined after the exploratory cohort was
opened, and the confirmatory endpoint after that cohort's recall and FPR were known.
Hide-and-Seek was chosen as the better of two supervised baselines on the opened cohort, and
registered after SAFE's confirmatory result was known, on cohorts already opened for SAFE. Part of SAFE + P's gain is its higher test FPR (0.061 against 0.055). Timers and normalization assume a known task and limit; the policy
score needs white-box access. Test
cohorts are balanced by quota, so recall and FPR are per class, not at deployment prevalence.
Point-of-no-return labels are one labeler's video judgments,
checked on 12 videos; only 10 failures were decided.

\section{Conclusion}

When failure means timeout, and mostly slowness, a learned monitor can only warn earlier than
a timer at the same false-alarm rate, here and on a published benchmark. Elapsed time explains
how well monitors separate failures. On fresh preregistered cohorts, SAFE's hidden-state score
added 0.030 of the time limit to a per-task timer, below our margin, while a sequence detector
trained for early separation on the same features added 0.071, a warning that develops during
the rollout. Stage~2 tests a second policy and closed-loop retry.

\bibliographystyle{plainnat}
\bibliography{refs}

\end{document}
